%! AP Statistics — Unit 1, Lesson 1.4 — Homework ANSWER KEY
\documentclass[10pt]{article}
\usepackage{apstats-article}
\usepackage{apstats-key}

%=============================================================================
\begin{document}

\pageheader{Unit 1, Lesson 1.4}{Homework --- Answer Key}

\vspace{0.1in}

\begin{tcolorbox}[colback=sky, colframe=navy, boxrule=0.8pt, arc=2mm,
    left=3mm, right=3mm, top=1.5mm, bottom=1.5mm]
\small For all problems: include a title and labeled axes on every graph; always
start the y-axis at 0; use complete sentences for written responses with context
included. AP free-response requires all of these --- build the habit now.
\end{tcolorbox}

\vspace{0.14in}

% ════════════════════════════════════════════════════════════════════════════
% PART A
% ════════════════════════════════════════════════════════════════════════════
\boxguard
\begin{blurbbox}[Part A \quad Construct a Bar Chart from a Frequency Table]{navylight}
\small

A survey of \textbf{40 middle school students} asks: ``How do you get to school?''
The results are below.

\vspace{8pt}
\renewcommand{\arraystretch}{1.8}
\begin{tabularx}{0.75\linewidth}{@{} >{\bfseries}p{0.22\linewidth}
                                        C{0.18\linewidth}
                                        C{0.26\linewidth} @{}}
\rowcolor{sky}
\textbf{Mode} & \textbf{Frequency} & \textbf{Rel.\ Frequency} \\
\midrule
Walk   & 16 & \ans{0.40 (40\%)} \\
Car    & 10 & \ans{0.25 (25\%)} \\
Bus    & 10 & \ans{0.25 (25\%)} \\
Bike   & 4  & \ans{0.10 (10\%)} \\
\midrule
\rowcolor{sky}\textbf{Total} & \ans{40} & \ans{1.00} \\
\end{tabularx}

\vspace{12pt}

\noindent\textbf{Construct a relative frequency bar chart} in the space below.

\vspace{8pt}
\begin{tcolorbox}[colback=white, colframe=slate, boxrule=0.5pt, arc=2mm,
    left=2mm, right=2mm, top=2mm, bottom=2mm, height=2.4in]
\centering
{\color{keyred}\bfseries\small How 40 Middle School Students Get to School}\\[2pt]
\begin{tikzpicture}[scale=0.92]
  % y-axis 0 to 0.45; 1 unit = 10cm
  \draw[thick] (0,0) -- (0,4.6);
  \foreach \p/\lbl in {0/0, 1/0.10, 2/0.20, 3/0.30, 4/0.40} {
    \draw (-0.1,\p) -- (0.1,\p);
    \node[left, font=\scriptsize] at (-0.12,\p) {\lbl};
  }
  \node[rotate=90, font=\scriptsize, color=keyred] at (-1.05,2.3) {Relative frequency};
  \draw[thick] (0,0) -- (7.0,0);
  \foreach \i/\v/\lbl in {0/4.0/Walk, 1/2.5/Car, 2/2.5/Bus, 3/1.0/Bike} {
    \fill[keyred!75] ({0.7+\i*1.6},0) rectangle ({1.7+\i*1.6},\v);
    \draw[keyred]    ({0.7+\i*1.6},0) rectangle ({1.7+\i*1.6},\v);
    \node[below, font=\scriptsize, color=keyred] at ({1.2+\i*1.6},-0.08) {\lbl};
  }
  \node[below, font=\scriptsize\itshape, color=keyred] at (3.5,-0.55) {Mode of transportation};
\end{tikzpicture}
\end{tcolorbox}

\vspace{12pt}

\begin{enumerate}[leftmargin=1.5em, itemsep=8pt]
  \item Write 2--3 complete sentences describing the distribution of
        transportation mode for these 40 students. Include the most and least
        common categories and their relative frequencies.\\[8pt]
        \ansline{Of the 40 middle school students surveyed, walking was the most}\\[1pt]
        \ansline{common mode (40\%). Car and bus tied at 25\% each. Biking was}\\[1pt]
        \ansline{least common, at only 10\%.}

  \item What percentage of students do not walk to school? Show your work.\\[8pt]
        \ansline{$100\% - 40\% = \textbf{60\%}$}\\[1pt]
        \ansline{(equivalently $(10+10+4)/40 = 24/40 = 0.60$).}

  \item Would a pie chart or bar chart be more appropriate here? Explain your
        reasoning.\\[6pt]
        \ansline{A bar chart --- comparing heights makes the Car/Bus tie at 25\%}\\[1pt]
        \ansline{obvious, where equal pie slices are hard to tell apart.}
\end{enumerate}

\end{blurbbox}

\vspace{0.14in}

% ════════════════════════════════════════════════════════════════════════════
% PART B
% ════════════════════════════════════════════════════════════════════════════
\boxguard
\begin{blurbbox}[Part B \quad Critique a Misleading Graph]{greenacc}
\small

A local newspaper publishes a bar chart showing the number of students enrolled
in each of four elective classes. The chart has the following properties:
\begin{itemize}[itemsep=3pt]
  \item The y-axis starts at 40 (not 0).
  \item The bar for Art class reaches 60; the bar for Music reaches 50.
  \item The headline reads: ``Art Has 20\% More Students Than Music!''
\end{itemize}

\vspace{10pt}

\begin{enumerate}[leftmargin=1.5em, itemsep=8pt]
  \item Explain why the newspaper's bar chart is misleading. What does a reader
        incorrectly perceive?\\[6pt]
        \ansline{The y-axis starts at 40, so the visible bars are $60-40=20$ against}\\[1pt]
        \ansline{$50-40=10$ --- Art looks twice Music. A reader perceives a 100\%}\\[1pt]
        \ansline{gap when the real difference is only 20\%.}

  \item Calculate the \textit{actual} percentage by which Art exceeds Music,
        relative to the true enrollment numbers. Is the headline accurate?\\[8pt]
        \ansline{$(60-50)/50 \times 100\% = \textbf{20\%}$. The headline is numerically}\\[1pt]
        \ansline{correct, but the graph makes 20\% look far larger than it is.}

  \item Describe exactly how you would fix the graph so it accurately represents
        the data.\\[6pt]
        \ansline{Redraw with the y-axis starting at 0 and label it with units.}\\[1pt]
        \ansline{Both bars then read near-equal, with Art modestly taller.}

  \item If the school has 300 total students enrolled in electives, and you
        know Art = 60 and Music = 50, what are their relative frequencies?
        Which class has the larger share of total enrollment?\\[6pt]
        \ansline{Art $= 60/300 = 0.200$ (20\%); Music $= 50/300 \approx 0.167$ (16.7\%).}\\[1pt]
        \ansline{Art has the larger share; the other 190 students take other electives.}
\end{enumerate}

\end{blurbbox}

\vspace{0.14in}

% ════════════════════════════════════════════════════════════════════════════
% PART C
% ════════════════════════════════════════════════════════════════════════════
\boxguard
\begin{blurbbox}[Part C \quad AP-Style Free Response]{redacc}
\small

The table below shows the distribution of color preference for a sample of
120 car buyers.

\vspace{8pt}
\renewcommand{\arraystretch}{1.6}
\begin{tabularx}{0.72\linewidth}{@{} >{\bfseries}p{0.28\linewidth}
                                      C{0.22\linewidth}
                                      C{0.26\linewidth} @{}}
\rowcolor{sky}
\textbf{Color} & \textbf{Frequency} & \textbf{Rel.\ Frequency} \\
\midrule
White   & 42 & 0.350 \\
Black   & 30 & 0.250 \\
Silver  & 24 & 0.200 \\
Red     & 12 & 0.100 \\
Other   & 12 & 0.100 \\
\midrule
\rowcolor{sky}\textbf{Total} & 120 & 1.000 \\
\end{tabularx}

\vspace{10pt}

\begin{enumerate}[leftmargin=1.5em, itemsep=10pt]
  \item Describe the distribution of car color preference for these 120 buyers.
        Use the information in the table in your response.\\[6pt]
        \ansline{Of the 120 car buyers, white was most preferred (35\%), then black}\\[1pt]
        \ansline{(25\%) and silver (20\%); red and other tied last at 10\% each.}\\[1pt]
        \ansline{White and black together account for 60\% of all preferences.}

  \item Compute the sector angle (in degrees) for each color in a pie chart.
        Which two colors together would form almost exactly half the circle?\\[8pt]
        \ansline{White $126^\circ$, Black $90^\circ$, Silver $72^\circ$, Red $36^\circ$, Other $36^\circ$.}\\[1pt]
        \ansline{Black $+$ Silver $= 90 + 72 = 162^\circ$, nearly half the $360^\circ$ circle}\\[1pt]
        \ansline{(together 45\% of preferences).}

  \item A car dealership wants to display this data on a poster. Should they
        use a bar chart or pie chart? Justify your choice in one or two sentences.\\[6pt]
        \ansline{A bar chart --- with 5 categories, heights compare more easily than}\\[1pt]
        \ansline{angles, and the equal Red/Other slices would be hard to tell apart.}
\end{enumerate}

\end{blurbbox}

\vspace{0.14in}

% ════════════════════════════════════════════════════════════════════════════
% EXTENSION
% ════════════════════════════════════════════════════════════════════════════
\boxguard
\begin{extensionbox}
\small
\textbf{Optional Extension --- Find a Misleading Graph in the Wild}\\[4pt]
Find a real bar chart or pie chart from a news article, social media, or a
company report that you believe is misleading. Write a 3--5 sentence analysis:
(1) identify the source and what data the graph displays, (2) explain specifically
why the graph is misleading, (3) describe how to correct the display, and (4) state
what a reader would wrongly conclude from the original graph.

\vspace{8pt}
\writeline\\[1pt]\writeline\\[1pt]\writeline\\[1pt]\writeline\\[1pt]\writeline

\vspace{4pt}
\textcolor{slate}{\scriptsize Source / URL:\;\blank{11cm}}
\end{extensionbox}

\end{document}
